\begin{table}[htbp]
\caption{Liberia Treatment Effects by Track and Within-Track Baseline Ability Tercile}
\label{tab:track_bins}
\begin{center}
\begin{small}
\begin{threeparttable}
\begin{tabular}{llcccc}
\toprule
Track & \Longstack{Ability\\ tercile} & Estimate & \Longstack{Standard\\ error} & $p$-value & $N$ \\
\midrule
Lower & Bottom & -0.165 & (0.205) & 0.423 & 460 \\
 & Middle & 0.208 & (0.266) & 0.435 & 459 \\
 & Top & -0.068 & (0.226) & 0.764 & 393 \\
Upper & Bottom & -0.198 & (0.172) & 0.249 & 626 \\
 & Middle & -0.484* & (0.255) & 0.057 & 606 \\
 & Top & -0.453*** & (0.173) & 0.009 & 610 \\
\bottomrule
\end{tabular}
\begin{tablenotes}
\footnotesize \item Notes: The table reports Liberia student-level intent-to-treat estimates for standardized endline reading scores, separately by assigned reading track and tercile of baseline ability within that track. Estimates come from OLS regressions with strata fixed effects and the empirical Bayes baseline-score control used in the main specifications. Standard errors are cluster-robust at the school-grade-group level. If treatment harm were driven mainly by marginal students being placed into an overly advanced upper track, the largest losses would appear in the bottom tercile of the upper track. Instead, the largest losses occur in the middle and top terciles of the upper track. $^*$ and $^{***}$ denote significance at the 10 and 1 percent levels, respectively.
\end{tablenotes}
\end{threeparttable}
\end{small}
\end{center}
\end{table}